% 第19章 t-J模型的半经典理论
\chapter{$t$--$J$ 模型的半经典理论}

层状铜氧化物（如 $La_{2-x}Sr_{x}CuO_{4}$）中高温超导电性的发现，激发了对近半满准二维 Hubbard 模型的深入研究。基本问题是理解反铁磁 Mott 型绝缘体（如 $La_{2}CuO_{4}$）如何在空穴浓度 $x=5$--$20\%$ 时演化为金属或超导体。

如 3.2 节所示，大 $U$ 时 Hubbard 模型的低能激发可以用 t--J 哈密顿量 (3.29) 描述：

\begin{equation*}
\mathcal {H} ^ {t - J} = P _ {s} \left( \mathcal {T} + \mathcal {H} ^ {\mathrm{QHM}} + \mathcal {J} ' \right) P _ {s} ,
\tag{19.1}
\end{equation*}

$P_{s}$ 把 Hilbert 空间投影到每格点零或一个电子的子空间，而

\begin{align*}
\mathcal {T} & = - t \sum_{\langle ij \rangle s} \left( c _ {is} ^ {\dagger} c _ {js} + c _ {js} ^ {\dagger} c _ {is} \right) ,\\
\mathcal {H} ^ {\mathrm{QHM}} & = J \sum_{\langle ij \rangle} \left( \mathbf {S} _ {i} \cdot \mathbf {S} _ {j} - \frac {n _ {i} n _ {j}}{4} \right) ,\\
\mathcal {J} ' & = - \frac {1}{2 U} \sum_{\langle i, j, k \rangle} ^ {i \neq k} t _ {ij} t _ {jk} \left( \sum_ {s} \left( c _ {is} ^ {\dagger} c _ {ks} \rho_ {j} \right) - c _ {i} ^ {\dagger} \vec {\sigma} c _ {k} \cdot c _ {j} ^ {\dagger} \vec {\sigma} c _ {j} \right) .
\tag{19.2}
\end{align*}

$\langle ij\rangle$ 是近邻键；$\langle i,j,k\rangle$ 表示对构成近邻三元组的格点 $i,j,k$ 求和。半满时 $n_{i}\to1$，跳跃项 $\mathcal{T}$ 与 $\mathcal{J}'$ 在投影子空间中消失，t--J 模型约化为自旋 1/2 的量子 Heisenberg 反铁磁体。

偏离半满时，t--J 模型描述相互作用自旋与可动空穴的体系，也称“掺杂反铁磁体”。

由于大 $U/t$ 蕴含 $t \gg J$，$t$ 跳跃项并不小。而且空穴在相反子晶格之间的跳跃会扰乱反铁磁关联。这带来一个根本性的理论困难：我们认为已理解得相当透彻的 Heisenberg 模型，被可动空穴的加入强烈微扰。这在极端的 Nagaoka 极限\footnote{见定理 4.1。}$J = 0$ 中表现得最为显著：单个空穴的基态是完全铁磁的！

为 t--J 模型选择近似方案时，最好能找到一个控制参数来估计误差。既然我们关心低掺杂浓度，而对 Heisenberg 极限的认识相当可靠，一个显然的小参数是空穴数目。另一个小参数是逆自旋大小 $1/S$，它控制半经典展开。

本章把半经典分析从 Heisenberg 模型扩展到少量空穴的 t--J 模型。先把 (19.1) 从自旋 1/2 推广到任意自旋 $S$，导出自旋--空穴相干态路径积分表示，可用最陡下降法以 $S$ 为大参数展开。发现经典解是在物理可及的 $t/J$ 区间内的小铁磁极化子。小极化子的半经典量子化导出低能自旋与电荷激发（有序相中）的有效 $t^{\prime}$--$J$ 哈密顿量。文中还提到自旋液体相中高温超导电性的一种可能方案。

\section{Schwinger 玻色子与从属费米子}

投影算符 $P_{s}$ 是电子态上的非完整（nonholonomic）约束：

\begin{equation*}
P _ {s} \left| n _ {i} \right\rangle = \left\{ \begin{array}{ll} \left| n _ {i} \right\rangle & n _ {i} \leq 1 \\ 0 & n _ {i} = 2 \end{array} \right. .
\tag{19.3}
\end{equation*}

通常，处理作为等式（而非不等式）的完整（holonomic）约束更容易。做法是在每个格点 $i$ 引入两个对易的 Schwinger 玻色子 $(a_{i}, b_{i})$ 和一个“从属费米子”$f_{i}$，满足标准代数

\begin{align*}
\left[ a _ {i}, a _ {j} ^ {\dagger} \right] & = \delta_ {ij} ,\\
\left[ b _ {i}, b _ {j} ^ {\dagger} \right] & = \delta_ {ij} ,\\
\left\{ f _ {i}, f _ {j} ^ {\dagger} \right\} & = \delta_ {ij} .
\tag{19.4}
\end{align*}

单占据空间中的电子态表示为

\begin{equation*}
\left\{ c _ {i \uparrow} ^ {\dagger} \left| 0 \right\rangle , c _ {i \downarrow} ^ {\dagger} \left| 0 \right\rangle , \left| 0 \right\rangle \right\} = \left\{ a _ {i} ^ {\dagger} \left| 0 \right\rangle , b _ {i} ^ {\dagger} \left| 0 \right\rangle , f _ {i} ^ {\dagger} \left| 0 \right\rangle \right\} ,
\tag{19.5}
\end{equation*}

$|0\rangle$ 是所有 Schwinger 玻色子与从属费米子的真空。$P_{s}$ 成为 Schwinger 玻色子与从属费米子乘积 Fock 空间上的完整约束：

\begin{equation*}
P _ {s} \left( a _ {i} ^ {\dagger} a _ {i} + b _ {i} ^ {\dagger} b _ {i} + f _ {i} ^ {\dagger} f _ {i} - 1 \right) = 0 \, , \quad \forall i.
\tag{19.6}
\end{equation*}

投影后的电子算符可表示为

\begin{align*}
c _ {i \uparrow} P _ {s} & \rightarrow f _ {i} ^ {\dagger} a _ {i} P _ {s} ,\\
c _ {i \downarrow} P _ {s} & \rightarrow f _ {i} ^ {\dagger} b _ {i} P _ {s} .
\tag{19.7}
\end{align*}

把 (19.7) 代入 (3.33) 的正规排序 t--J 模型：

\begin{align*}
\mathcal {H} ^ {t - J} & = P _ {s} \Bigg[ t \sum_{\langle ij \rangle} \left( f _ {i} ^ {\dagger} f _ {j} \mathcal {F}_{ji} + f _ {j} ^ {\dagger} f _ {i} \mathcal {F}_{ij} \right)\\
& \quad - \frac {J}{4} \sum_{\langle i, j, k \rangle} \left( \delta_{ik} - f _ {k} ^ {\dagger} f _ {i} \right) \mathcal {A}_{ij} ^ {\dagger} \mathcal {A}_{kj} \left( 1 - f _ {j} ^ {\dagger} f _ {j} \right) \Bigg] P _ {s} ,
\tag{19.8}
\end{align*}

$\langle ij\rangle$ 表示近邻键，$\langle i,j,k\rangle$ 表示三个相邻格点的序列。键算符按 (16.2) 与 (16.10) 定义：

\begin{align*}
\mathcal {F} _ {ij} & = a _ {i} ^ {\dagger} a _ {j} + b _ {i} ^ {\dagger} b _ {j} ,\\
\mathcal {A} _ {ij} & = a _ {i} b _ {j} - b _ {i} a _ {j} .
\tag{19.9}
\end{align*}

半满（未掺杂）时 $n_{i}^{f} = 0$，Schwinger 玻色子代表纯自旋 1/2 算符；(19.8) 中跳跃项消失，第二项约化为反铁磁 Heisenberg 模型（见 (16.13)）。

用这一表示，只需把约束 (19.6) 换成

\begin{equation*}
P _ {S} \left( a _ {i} ^ {\dagger} a _ {i} + b _ {i} ^ {\dagger} b _ {i} + f _ {i} ^ {\dagger} f _ {i} - 2S \right) = 0 \, , \quad \forall i.
\tag{19.10}
\end{equation*}

就能把 t--J 模型从自旋 1/2 推广到任意自旋 $S$。这一约束意味着：格点 $i$ 上有空洞时，该处自旋减小 $\frac{1}{2}$。顺带一提，大 $S$ 自旋--空穴表示可能在过渡金属离子中实现\footnote{若干 $d$ 电子按 Hund 定则铁磁耦合的情形；见 2.1 节。}。

\section{自旋--空穴相干态}

我们希望构造满足约束 (19.10) 的自旋 $S$ t--J 模型 (19.8) 的表示。先回忆 7.3 节的自旋相干态：

\begin{equation*}
\left| \hat {\Omega} \right\rangle_{S} \equiv e ^ {iS\chi} \left[ ( 2S )! \right] ^ {-1/2} \left[ u ( \theta, \phi ) a ^ {\dagger} + v ( \theta, \phi ) b ^ {\dagger} \right] ^ {2S} \left| 0, 0 \right\rangle .
\tag{19.11}
\end{equation*}

$\chi$ 是任意相位，$\hat{\Omega}$ 是纬度 $\theta$、经度 $\phi$ 的单位矢量，而

\begin{align*}
u & = \cos ( \theta / 2 ) \, e ^ {i\phi/2} ,\\
v & = \sin ( \theta / 2 ) \, e ^ {-i\phi/2} .
\tag{19.12}
\end{align*}

自旋相干态在自旋 $S$ 子空间提供单位分解：

\begin{equation*}
\frac {2S + 1}{4 \pi} \int d \hat {\Omega} \left| \hat {\Omega} \right\rangle \left\langle \hat {\Omega} \right| = I ,
\tag{19.13}
\end{equation*}

$d\hat{\Omega}=d\phi\; d\cos\theta$。第 10 章曾用这一分解构造自旋相干态路径积分。这里把基底扩展到 (19.10) 的 Hilbert 空间，定义“自旋--空穴相干态”：

\begin{equation*}
\left| \hat {\Omega}, \xi \right\rangle_{S} \equiv \left| \hat {\Omega} \right\rangle_{S} \left| 0 \right\rangle_{f} + \left| \hat {\Omega} \right\rangle_{S - \frac {1}{2}} \xi f ^ {\dagger} \left| 0 \right\rangle_{f} ,
\tag{19.14}
\end{equation*}

$\xi$ 是反对易 Grassmann 变量\footnote{见附录 C。}。

两个态（$\hat{\Omega} \approx \hat{\Omega}'$）的交叠为

\begin{align*}
\left\langle \hat {\Omega}, \xi \middle| \hat {\Omega} ', \xi ' \right\rangle & = \langle \hat {\Omega} | \hat {\Omega} ' \rangle_{S} \left[ 1 + \xi ^ {*} \xi ' \left( \langle \hat {\Omega} | \hat {\Omega} ' \rangle_{\frac {1}{2}} \right) ^ {- 1} \right]\\
& \simeq \exp \left( - iS \psi [ \hat {\Omega}, \hat {\Omega} ' ] + e ^ {i \frac {1}{2} \psi [ \hat {\Omega}, \hat {\Omega} ' ]} \xi ^ {*} \xi ' \right) ,
\tag{19.15}
\end{align*}

其中

\begin{align*}
\psi & = 2 \arctan \left\{ \tan \left( \frac {\phi - \phi '}{2} \right) \frac {\cos \left[ \frac {1}{2} ( \theta + \theta ' ) \right]}{\cos \left[ \frac {1}{2} ( \theta - \theta ' ) \right]} \right\} + \chi - \chi '\\
& \simeq - \mathbf {A} \cdot \left( \hat {\Omega} - \hat {\Omega} ' \right) ,
\tag{19.16}
\end{align*}

$\mathbf{A}$ 是 (10.17) 或 (10.18) 给出的熟悉的磁单极矢量势，满足

\begin{equation*}
\nabla \times \mathbf {A} \cdot \hat {\Omega} = 1.
\tag{19.17}
\end{equation*}

(19.11) 中选择 $\chi(\hat{\Omega})$ 的自由等价于 $\mathbf{A}$ 的规范自由。于是由 (19.15)，$\mathbf{A}$ 的规范变换必须伴随 Grassmann 变量的 $U(1)$ 相位变换，反之亦然：

\begin{align*}
\chi ( \hat {\Omega} ) & \rightarrow \chi ' ( \hat {\Omega} ) ,\\
\mathbf {A} & \rightarrow \mathbf {A} ' = \mathbf {A} + \nabla_{\hat {\Omega}} \left( \chi ' - \chi \right) ,\\
\xi & \rightarrow e ^ {i ( \chi ' - \chi )/2} \xi .
\tag{19.18}
\end{align*}

这展示了 (19.15) 中规范协变耦合与局域约束 (19.10) 之间的密切关系。

自旋--空穴相干态在受约束 Hilbert 空间中提供单位分解：

\begin{equation*}
\frac {2S}{4 \pi} \int d \hat {\Omega} \, d \xi^ {*} \, d \xi \, \exp \left[ - \alpha_{S} \xi^ {*} \xi \right] \left| \hat {\Omega}, \xi \right\rangle \left\langle \hat {\Omega}, \xi \right| = I ,
\tag{19.19}
\end{equation*}

这里用了 (19.14) 与 Grassmann 的 Gauss 积分 (C.10)。Gauss 归一化因子为

\begin{equation*}
\alpha _ {S} = \frac {2S + 1}{2S}.
\tag{19.20}
\end{equation*}

现在可以在每个离散时间步 $\epsilon = \beta/N_{\epsilon}$ 之间插入单位分解，构造虚时生成泛函的路径积分表示。于是得到\footnote{见第 10 章与附录 D。}

\begin{align*}
Z & = \lim_{N _ {\epsilon} \to \infty} \int \mathcal {D} \hat {\Omega} \, \mathcal {D} ^ {2} \xi \, \exp \left[ - \sum_{i, \tau} \alpha_{S} \xi_{i} ^ {*} ( \tau ) \xi_{i} ( \tau ) \right]\\
& \quad \times \prod_{\tau = \epsilon} ^ {\beta} \left\langle \hat {\Omega} ( \tau ), \xi ( \tau ) \left| 1 - \epsilon \mathcal {H} \right| \hat {\Omega} ( \tau - \epsilon ), \xi ( \tau - \epsilon ) \right\rangle .
\tag{19.21}
\end{align*}

连续极限为

\begin{align*}
Z & = \int \mathcal {D} \hat {\Omega} \, \mathcal {D} ^ {2} \xi \, \exp \Big\{ \int_ {0} ^ {\beta} d \tau \sum_ {i} \left[ i \left( S - \frac {1}{2} \xi_{i} ^ {*} \xi_{i} \right) \mathbf {A} ( \hat {\Omega} _ {i} ) \cdot \dot {\hat {\Omega}} _ {i} - \xi_{i} ^ {*} \partial_ {\tau}^{(\alpha)} \xi_{i} \right]\\
& \quad \left. - H ^ {t - J} [ \hat {\Omega}, \xi ^ {*}, \xi ] + \mu \sum_ {i} \xi_{i} ^ {*} \xi_{i} \right\} .
\tag{19.22}
\end{align*}

回忆\footnote{自旋动能项的讨论见第 10、11 章。}把时间导数当作 $\hat{\Omega}(\tau)$ 可微来处理，是对算符排序的粗心，会给基态能量带来不确定性。对费米子，无需假设 $\xi(\tau)$ 连续。离散时间导数为

\begin{equation*}
\int d \tau \left( \xi_{i} ^ {*} \partial_ {\tau}^{(\alpha)} \xi_{i} \right) \equiv \sum_ {\tau} \xi_{i} ^ {*} ( \tau ) \left[ \alpha_{S} \xi_{i} ( \tau ) - \xi_{i} ( \tau - \epsilon ) \right].
\tag{19.23}
\end{equation*}

（不必担心 $\alpha_{S} \neq 1$，它很快会被固定。）

哈密顿量函数为

\begin{align*}
H ^ {t - J} & = \frac {\left\langle \hat {\Omega}, \xi \left| \mathcal {H} ^ {t - J} \right| \hat {\Omega}, \xi \right\rangle}{\langle \hat {\Omega}, \xi | \hat {\Omega}, \xi \rangle} = H ^ {t} + H ^ {J} ,\\
H ^ {t} & = \frac {\bar {t}}{\sqrt {2}} \sum_{\langle ij \rangle} \sqrt {1 + \hat {\Omega} _ {i} \cdot \hat {\Omega} _ {j}} \left( e ^ {\frac {i}{2} \psi_{ij} ^ {F}} \xi_{i} ^ {*} \xi_{j} + e ^ {- \frac {i}{2} \psi_{ij} ^ {F}} \xi_{j} ^ {*} \xi_{i} \right) ,\\
H ^ {J} & = - \frac {\bar {J}}{8} \sum_{\langle i, j, k \rangle} e ^ {\frac {i}{2} ( \psi_{ij} ^ {N} - \psi_{jk} ^ {N} )} \sqrt { \left( 1 - \hat {\Omega} _ {j} \cdot \hat {\Omega} _ {k} \right) \left( 1 - \hat {\Omega} _ {i} \cdot \hat {\Omega} _ {j} \right) }\\
& \quad \times \left( 1 - \xi_{j} ^ {*} \xi_{j} \right) \xi_{i} ^ {*} \xi_{k} ,
\tag{19.24}
\end{align*}

（依赖规范的）相位 $\psi^F, \psi^N$ 已在 (17.5) 与 (17.10) 定义。

为了在 $S$ 大时给 $H^{t-J}$ 一个有限的经典极限，把 (19.24) 的 $S$ 依赖吸收进如下标度的经典参数：

\begin{align*}
4JS ^ {2} & \to \bar {J} ,\\
2tS & \to \bar {t} .
\tag{19.25}
\end{align*}

（这些定义在 $S = \frac{1}{2}$ 时还原为原参数 $J, t$。）

$H^{J'}$ 含 Grassmann 变量的四次幂，难以积分，需要近似。在 Fock 基中，四次项的大小为

\begin{equation*}
\left| \langle f _ {i} ^ {\dagger} f _ {k} f _ {j} ^ {\dagger} f _ {j} \rangle_{\mathrm{Fock}} \right| = \left| \rho_{ik} \right| \rho_{jj} - \delta_{ik} \rho_{ij} \rho_{ji} ,
\tag{19.26}
\end{equation*}

$\rho$ 是空穴密度矩阵。我们的经典解中，空穴密度主要位于 $\sqrt{1-\hat{\Omega}_{i}\cdot\hat{\Omega}_{j}}\approx0$ 的铁磁区。由 (19.24)，这些正是 $\rho_{ij}$ 特别小的地方。因此可以放心地把 $H^{t-J}$ 换成简化的双线性哈密顿量\footnote{哈密顿量 (19.27) 在文献中常被直接称作 t--J 模型。}：

\begin{align*}
H ^ {t - J} & \approx \sum_ {ij} \hat {H}_{ij} ^ {t} \xi_{i} ^ {*} \xi_{j} + \bar {H} ^ {J} [ \hat {\Omega}, \rho ] ,\\
\hat {H} _ {ij} ^ {t} & = \frac {\bar {t}}{\sqrt {2}} \sqrt {1 + \hat {\Omega} _ {i} \cdot \hat {\Omega} _ {j}} \; e ^ {i \psi_{ij} ^ {F}} ,\\
\bar {H} ^ {J} & = \frac {\bar {J}}{4} \sum_{\langle ij \rangle} \left( \hat {\Omega} _ {i} \cdot \hat {\Omega} _ {j} - 1 \right) \left( 1 - \rho_{i} \right) \left( 1 - \rho_{j} \right) .
\tag{19.27}
\end{align*}

第三行把三元组 $\langle i,j,k\rangle$ 的求和换成键 $\langle ij\rangle$ 的求和，并略去非对角跳跃项（与 $t$ 跳跃项相比是小量）。

现在积掉 Grassmann 变量\footnote{见附录 D 的 D.2 节。}，得到纯自旋路径积分：

\begin{align*}
Z & = \int \mathcal {D} \hat {\Omega} \, \exp \Big\{ \int_ {0} ^ {\beta} d \tau \Big[ i \sum_ {i} \left( S - \rho_{i} /2 \right) \mathbf {A} ( \hat {\Omega} _ {i} ) \cdot \dot {\hat {\Omega}} _ {i} - \bar {H} ^ {J} [ \hat {\Omega} ]\\
& \quad \left. - F ^ {f} [ \hat {\Omega}, \mu ' ] \Big] \right\} .
\tag{19.28}
\end{align*}

$F^{f}[\mu]$ 是空穴的（推迟）自由能泛函：

\begin{align*}
F ^ {f} & = - \beta^ {- 1} \operatorname{Tr}_{i} \ln \left[ \alpha_{S} I + T_{\tau} \exp \left[ - \int_ {0} ^ {\beta} d \tau \left( \hat {H}_{ij} ^ {t} [ \hat {\Omega} ( \tau ) ] - \mu \delta_{ij} \right) \right] \right]\\
& = - \beta^ {- 1} \operatorname{Tr}_{i} \ln \left[ I + T_{\tau} \exp \left[ - \int_ {0} ^ {\beta} d \tau \left( \hat {H}_{ij} ^ {t} [ \hat {\Omega} ( \tau ) ] - \mu ' \delta_{ij} \right) \right] \right]\\
& \quad - \frac {1}{\epsilon} \mathcal {N} \ln \left( \alpha_{S} \right) .
\tag{19.29}
\end{align*}

第二行把一个无穷大常数吸收化学势：

\begin{equation*}
\mu ' = \mu - \frac {1}{\epsilon} \ln \alpha_{S}.
\tag{19.30}
\end{equation*}

化学势与基态能量的平移对关联函数无关紧要；此后忽略之，取 $\alpha_{S} \rightarrow 1$ 与 $\mu' \rightarrow \mu$。

方程 (19.28) 是半经典近似的有用出发点。$S \to \infty$ 时动能项系数增大，压低含时自旋涨落。这使我们能把 $F^{f}[\hat{\Omega}]$ 与密度 $\rho[\hat{\Omega}]$ 视为 $\hat{\Omega}(\tau)$ 的绝热函数。也就是说：大 $S$ 同时证明了半经典近似与绝热近似。

\section{经典理论：小极化子}

严格的经典极限下，路径积分完全由不依赖时间的组态主导：

\begin{equation*}
\hat {\Omega} ( \tau ) \rightarrow \hat {\Omega} ( 0 ) = \hat {\Omega} ( \beta ).
\tag{19.31}
\end{equation*}

由 (10.23)，$Z$ 正比于一个经典配分函数：

\begin{equation*}
Z \rightarrow Z ' \int \prod_ {i} d \hat {\Omega} _ {i} \, \exp \left[ - \beta \left( \bar {H} ^ {J} [ \hat {\Omega} ] + F ^ {f} [ \hat {\Omega} ] \right) \right].
\tag{19.32}
\end{equation*}

零温时，固定空穴数（$N^{f}$）的经典基态为

\begin{align*}
E ^ {cl} [ \hat {\Omega} ] & = \bar {H} ^ {J} [ \hat {\Omega} ] + \lim_{T \to 0} F ^ {f} [ \hat {\Omega} ] ,\\
E _ {0} ^ {cl} & = \left. \min_{\hat {\Omega}} E ^ {cl} \right| _ {\hat {\Omega} ^ {cl}, \, N ^ {f}} ,
\tag{19.33}
\end{align*}

先考察单个空穴的情形。

对单空穴极小化 $E_{0}^{cl}$ 概念上直截了当。物理上有两种相互竞争的相互作用：$F^{f}$ 偏爱尽可能多的高度连通的铁磁键，以使空穴动能最小；$\bar{H}^{J}$ 当然偏爱反铁磁关联。候选经典组态或是：

\begin{enumerate}[itemsep=1pt]
\item 弱畸变（倾倒或螺旋）的 Néel 组态，空穴密度铺展在整个格子上；或
\item 局域缺陷：空穴密度与自旋畸变集中在某个特定格子位置附近。
\end{enumerate}

变分极小化\footnote{见 Auerbach 与 Larson。}发现：当

\begin{equation*}
1 < \bar {t} / \bar {J} < 4
\tag{19.34}
\end{equation*}

时经典组态属第二类：五格点极化子，画于图~\ref{fig:19.1}。大 $\bar{t} / \bar{J} \gg 1$ 时极化子半径按\footnote{见习题 1。}

\begin{equation*}
R _ {c} \sim \left( \bar {t} / \bar {J} \right) ^ {1/4}
\tag{19.35}
\end{equation*}

增长。$\bar{t}/\bar{J} = \infty$ 极限下整个格子变成完全极化的铁磁体。有趣的是，这与无穷大 $U$ Hubbard 模型的 Nagaoka 定理 4.1 精确结果一致。我们还学到：对单空穴达到 Nagaoka 极限需要 $U/t$ 超过 $\mathcal{O}(\mathcal{N}^{2})$。

五格点极化子不扭曲背景组态的反铁磁关联（不像比如单条铁磁键会在远处诱导偶极形状的畸变）。极化子边界的“锐利”可解释如下：空穴能与 Heisenberg 能都对极化子边界处的畸变二次依赖；较强的能量胜出，故边界自旋与 Néel 方向完全对齐。这意味着超出短距离后，两个极化子之间没有经典相互作用。

\begin{figure}[H]
\centering
\includegraphics[width=0.45\textwidth]{fig19_1.jpg}
\caption{五格点极化子。圆圈描绘空穴密度。}
\label{fig:19.1}
\end{figure}

空穴的费米子谱有约 $2\bar{t}$ 的能隙。典型的自旋频率为

\begin{equation*}
\omega \sim \mathcal {O} ( \bar {J} S ^ {- 1} ) \ll \mathcal {O} ( \bar {t} ) ,
\tag{19.36}
\end{equation*}

因此空穴自由能的绝热近似是合理的。这与 Born--Oppenheimer 近似类似：慢自旋像重原子核，空穴像快电子，为慢变量确定一个有效势。

以格点 $i$ 为心的五格点极化子的自旋--空穴波函数为

\begin{equation*}
\left| \hat {\Omega} ^ {i}, \rho \right\rangle = \sum_{j = i, \, \langle ij \rangle} \sqrt {\rho_{j} [ \hat {\Omega} ^ {i} ]} \left| \hat {\Omega}_{j} ^ {i} \right\rangle_{S - \frac {1}{2}} f _ {j} ^ {\dagger} \left| 0 \right\rangle_{f _ {j}} \prod_{j ' \neq j} \left( \left| \hat {\Omega}_{j '} ^ {i} \right\rangle_{S} \left| 0 \right\rangle_{f _ {j '}} \right) ,
\tag{19.37}
\end{equation*}

$\rho[\hat{\Omega}^i]$ 是极化子组态 $\hat{\Omega}^i$ 的基态空穴密度。$S = \frac{1}{2}$ 时，(19.37) 可从纯 Néel 态 $\hat{\Omega}^N$ 通过

\begin{equation*}
\left| \hat {\Omega} ^ {i}, \rho \right\rangle = p _ {is _ {i}} ^ {\dagger} \left| \hat {\Omega} ^ {N} \right\rangle
\tag{19.38}
\end{equation*}

产生，$p^{\dagger}$ 是 Dagotto 与 Schrieffer 的准粒子算符

\begin{equation*}
p _ {is _ {i}} ^ {\dagger} \equiv \left( \sqrt {\rho_{i} [ \hat {\Omega} ^ {i} ]} + \sum_{\langle ij \rangle, s '} \sqrt {\rho_{j} [ \hat {\Omega} ^ {i} ]} \, c_{is '} ^ {\dagger} c_{js '} \right) c_{is}.
\tag{19.39}
\end{equation*}

$s_i = \pm \frac{1}{2}$ 是 $\hat{\Omega}_i^N$ 的自旋方向。

\section{极化子动力学与自旋隧穿}

五格点极化子破缺哈密顿量的自旋转动对称与格子平移对称。有限 $S$ 下，即使在零温，对称性也可由自旋的量子涨落恢复。恢复对称的涨落分两类：

\begin{enumerate}[itemsep=1pt]
\item 自旋波：倾向于恢复自旋对称，贡献在 $\mathcal{O}(1/S)$ 阶；
\item 自旋隧穿效应：移动极化子中心，恢复格子平移对称，是 $\mathcal{O}(e^{-S(\cdot)})$ 阶的非微扰效应。
\end{enumerate}

自旋波动力学可分别用第 11 章或第 12 章的方法处理有序相与无序相。

极化子的平移对称是分立的，自旋的任何运动都要爬越能垒。不过格子对称可以通过隧穿事件恢复：极化子的中心自旋转动到与近邻反铁磁关联，同时另一处的自旋转动形成新的五格点铁磁极化子（见图~\ref{fig:19.2}）。极化子从格点 $i$ 运动到 $k$ 的半经典隧穿矩阵元为\footnote{正确的多维隧穿矩阵元涉及受限制的波函数与磁通算符（见 Auerbach 与 Kivelson）。它们给出正确的前置因子 $\Gamma_{0}$，而 (19.40) 可用于领头阶行为。}

\begin{align*}
\Gamma_{ik} & \propto \left\langle \hat {\Omega} ^ {i} \left| G \left( E _ {0} ^ {cl} \right) \right| \hat {\Omega} ^ {k} \right\rangle\\
& = \int_{\hat {\Omega} ^ {i}} ^ {\hat {\Omega} ^ {k}} \mathcal {D} \hat {\Omega} ( t ) \, \exp \Bigg\{ i \int_ {0} ^ {\infty} d t \Big[ \sum_ {i} \left( S - \rho_{i} /2 \right) \left( 1 - \cos \theta_ {i} \right) \dot {\phi} _ {i}\\
& \quad \left. - \left( \bar {H} ^ {J} + E ^ {f} - E _ {0} ^ {cl} \right) \Big] \right\} ,
\tag{19.40}
\end{align*}

$G(E)$ 是 (10.38) 定义的依赖能量的 Green 函数。由转动对称，$H[\hat{\Omega}]$ 在所有方位角的统一平移

\begin{equation*}
\phi _ {i} = \bar {\phi} + \phi _ {i} '
\tag{19.41}
\end{equation*}

下不变。对动能项分部积分：

\begin{align*}
& \int \mathcal {D} \bar {\phi} ( t ) \exp \left[ - i \int d t \, \bar {\phi} ( t ) \, \partial_ {t} \left( \sum_ {i} \left[ S - \rho_{i} /2 ( t ) \right] \left[ 1 - \cos \theta_ {i} ( t ) \right] \right) \right]\\
& \propto \delta \left( M ^ {z} [ \hat {\Omega} ^ {i} ] - M ^ {z} [ \hat {\Omega} ^ {k} ] \right) ,
\tag{19.42}
\end{align*}

其中

\begin{equation*}
M ^ {z} [ \hat {\Omega} ( t ) ] = \sum_ {j} \left[ S - \rho_{j} ( t ) /2 \right] \left[ 1 - \cos \theta_ {j} ( t ) \right] = M ^ {z} [ \hat {\Omega} ( 0 ) ].
\tag{19.43}
\end{equation*}

也就是说，$z$ 磁化对 $G(E)$ 的任何矩阵元守恒。当然，这是哈密顿量转动不变性的精确推论。

一个重要结论是：在长程有序相中，五格点极化子不能在子晶格 A、B 之间隧穿，因为那会改变总磁化。换言之，存在一个禁止极化子跨子晶格跳跃的选择定则。这一强结果有些出人意料，因为原 t--J 哈密顿量 (19.27) 有大的跨子晶格跳跃项。物理上，原来的 $f$ 空穴被牢牢束缚于局域自旋缺陷，把低能下的电荷可动性限制在子晶格内的跳跃上。但要注意：若两个相距很远的格点所在处的局域环境具有反转的 Néel 场，极化子可以在分属相反子晶格的这两点之间隧穿（例如 Skyrmion 组态的中心与尾部）。

\begin{figure}[H]
\centering
\includegraphics[width=0.45\textwidth]{fig19_2.jpg}
\caption{极化子的隧穿路径。}
\label{fig:19.2}
\end{figure}

隧穿路径是 (19.40) 的鞍点，连接 $\hat{\Omega}^i\to \hat{\Omega}^k$ 并满足变分运动方程

\begin{equation*}
\frac {\delta}{\delta \hat {\Omega} _ {i} ( \tau )} \int_ {0} ^ {\infty} d \tau ' \left[ \left( S - \rho_{i} /2 \right) \left( 1 - \cos \theta_ {i} \right) \dot {\phi} _ {i} - E ^ {cl} [ \hat {\Omega} ] \right] = 0.
\tag{19.44}
\end{equation*}

（我们已固定 $\mathbf{A}$ 的规范选择来进行隧穿率计算。）

确定隧穿路径至少可以用两条守恒规则：

\begin{enumerate}[itemsep=1pt]
\item 沿经典路径能量守恒（见 (10.37)）：

\begin{equation*}
E ^ {cl} [ \hat {\Omega}^{\mathrm{tun}} ( \tau ) ] = E _ {0} ^ {cl}.
\tag{19.45}
\end{equation*}

\item 总磁化守恒 (19.43)。
\end{enumerate}

由于 $\hat{\Omega}^{\mathrm{tun}}$ 穿越经典禁戒区域，只有把坐标复化才能找到解\footnote{自旋路径积分类似相空间路径积分，共轭变量是 $\cos\theta, \phi$，见 (11.7) 与 (11.5)。允许 $\phi$ 变虚等价于让动量变虚，即把时间延拓到 Feynman 路径积分的虚轴。}。设自旋沿 $\pm\hat{y}$ 方向有序，用 $\{\hat{\Omega}_{i}^{\mathrm{tun}}(\varphi)\}$ 参数化隧穿路径，$\varphi$ 是自旋 $i$ 的方位转动角。

把隧穿路径的 $\phi$ 变量复化：

\begin{equation*}
\phi ( \varphi ) \rightarrow \phi^ {\mathrm{re}} ( \varphi ) + i \phi^ {\mathrm{im}} ( \varphi ).
\tag{19.46}
\end{equation*}

用 (19.45)，(19.40) 的领头阶最陡下降近似为

\begin{align*}
\Gamma_{ik} & \approx - \sum_{\alpha} \Gamma_{0}^{\alpha} \exp \left( - W_{ik}^{\alpha} + i \Upsilon_{ik}^{\alpha} \right) ,\\
W_{ik} & = \int d \varphi \sum_{j} \frac {d \phi_{j}^{\mathrm{im}}}{d \varphi} \left( S - \rho_{j} /2 \right) \left( 1 - \cos \theta_{j} \right) ,\\
\Upsilon_{ik}^{\alpha} & = \int d \varphi \sum_{j} \frac {d \phi_{j}^{\mathrm{re}}}{d \varphi} \left( S - \rho_{j} /2 \right) \left( 1 - \cos \theta_{j} \right) ,
\tag{19.47}
\end{align*}

$\alpha$ 标记隧穿路径；$\Gamma_{0}^{\alpha}$ 是正的前置因子；整体负号是时间反演不变哈密顿量中简并极小之间隧穿矩阵元的一般特征。

由 (19.47) 可得一些定性信息。第一，$|\Gamma_{ik}|$ 随 $S$ 指数减小。第二，$\Gamma$ 的符号由 $\Upsilon$——一个 Berry 相位——决定。现在来确定五格点极化子隧穿的 $\Upsilon$。

由 $\phi_{i}\rightarrow-\phi_{i}$ 下的对称性，存在两条隧穿路径 $\hat{\Omega}^{\mathrm{tun},+}$ 与 $\hat{\Omega}^{\mathrm{tun},-}$，分别对应 $xy$ 平面内的顺时针与逆时针转动：

\begin{equation*}
\phi_{i}^{\mathrm{re}, +} ( \varphi ) = - \phi_{i}^{\mathrm{re}, -} ( \varphi ).
\tag{19.48}
\end{equation*}

于是两条路径有相同的 $W_{ik}$ 但相反的相位 $\Upsilon$。对两条路径求和：

\begin{align*}
\Gamma_{ik} \approx & - \Gamma_{0} \exp \left[ - \int_ {0} ^ {\pi} d \varphi \sum_{i} \left( S - \frac {\rho_{i}}{2} \right) \left( 1 - \cos \theta_{i} \right) \frac {d \phi_{i}^{\mathrm{im}}}{d \varphi} \right]\\
& \times \sum_{\pm} \exp \left[ i \int_ {0} ^ {\pm \pi} d \varphi \sum_{i} \left( S - \frac {\rho_{i}}{2} \right) \left( 1 - \cos \theta_{i} \right) \frac {d \phi_{i}^{\mathrm{re}}}{d \varphi} \right]\\
= & - 2 e ^ {i 2 \pi S} \left| \Gamma_{ik} \right| ,
\tag{19.49}
\end{align*}

其中

\begin{align*}
\left| \Gamma_{ik} \right| & = - \Gamma_{0} \cos \left( \int_ {0} ^ {\pm \pi} d \varphi \sum_{i} \frac {\rho_{i}}{2} \frac {d \phi_{i}^{\mathrm{re}}}{d \varphi} \right)\\
& \quad \times \exp \left[ - \int_ {0} ^ {\pi} d \varphi \sum_{i} \left( S - \frac {\rho_{i}}{2} \right) \left( 1 - \cos \theta_{i} \right) \frac {d \phi_{i}^{\mathrm{im}}}{d \varphi} \right] .
\tag{19.50}
\end{align*}

利用 $\sum_{i}\rho_{i}=1$ 可确立因子 $\cos(\int d\varphi\ldots)$ 为正。相位 $2\pi S$ 来自 $\phi_{i}$ 与 $\phi_{k}$ 的整体转动；其他自旋折返原路径，其 Berry 相位或为零、或与对称伙伴干涉相消。

在 (19.49) 中我们发现：对半奇整数自旋，极化子跳跃的符号为正。这影响基态动量如下。

设完美的 Néel 背景。极化子跳跃的能带结构为

\begin{align*}
\epsilon_{\mathbf {k}} & = \sum_ {j} \Gamma_{ij} \exp \left[ i \mathbf {k} \left( \mathbf {x} _ {i} - \mathbf {x} _ {j} \right) \right]\\
& \approx 2 \Gamma_{(2, 0)} \left[ \cos \left( 2 k _ {x} \right) + \cos \left( 2 k _ {y} \right) \right]\\
& \quad + 2 \Gamma_{(1, 1)} \left[ \cos \left( k _ {x} + k _ {y} \right) + \cos \left( k _ {x} - k _ {y} \right) \right] ,
\tag{19.51}
\end{align*}

$(1,1)$、$(2,0)$ 表示同一子晶格上的第一、第二近邻格点（见图~\ref{fig:19.3}）。更远程的跳跃隧穿率更小，忽略之。由 (19.49)，整数自旋（$\Gamma < 0$）时 $\epsilon_{k}$ 在 $k = 0$ 极小。半奇整数自旋时，基态动量位于线

\begin{equation*}
\mathbf {k} _ {0} \in \left\{ \mathbf {k} : \cos k _ {x} + \cos k _ {y} = 0 \right\}.
\tag{19.52}
\end{equation*}

$\Gamma_{ij}$ 的符号之所以重要，是因为子晶格内跳跃涉及三角形（见图~\ref{fig:19.3}），其负号无法通过重定义波函数或在 (19.44) 中选不同规范约定而“规范掉”。

半经典理论解释了小团簇上自旋 1/2 t--J 模型的数值研究\footnote{见 Dagotto 与 Schrieffer。}。第一，基态动量位于 $(0,\pi)$ 或 $(\pi/2,\pi/2)$（及其对称反演），与 (19.52) 一致。第二，准粒子权重因子

\begin{figure}[H]
\centering
\includegraphics[width=0.45\textwidth]{fig19_3.jpg}
\caption{主导的极化子跳跃路径。}
\label{fig:19.3}
\end{figure}

\begin{equation*}
Z _ {h} = \left| \left\langle \Psi_ {0} ^ {t - J} \left| p_{\mathbf {k} _ {0} s} ^ {\dagger} \right| \Psi_ {0} ^ {\mathrm{Heis}} \right\rangle \right| ^ {2} / \left\langle \Psi_ {0} ^ {t - J} \left| p_{\mathbf {k} _ {0} s} p_{\mathbf {k} _ {0} s} ^ {\dagger} \right| \Psi_ {0} ^ {\mathrm{Heis}} \right\rangle
\tag{19.53}
\end{equation*}

可在小格子上测量。这里

\begin{equation*}
p _ {\mathbf {k} s} ^ {\dagger} = \sum_ {i} e ^ {- i \mathbf {k} \cdot \mathbf {x} _ {i}} p _ {is} ^ {\dagger}
\tag{19.54}
\end{equation*}

$p_{is}^{\dagger}$ 由 (19.39) 给出。$\Psi_{0}^{\mathrm{QHM}}, \Psi_{0}^{\mathrm{t-J}}$ 分别是 Heisenberg 模型与带一个空穴的 t--J 模型的精确基态。大 $S$ 极限下预期 $Z_{h} \rightarrow 1$。Dagotto 与 Schrieffer 数值计算了 $4 \times 4$ 正方格子 t--J 模型的 $Z_{h}$，在参数区 (19.34) 中发现

\begin{equation*}
0. 7 < Z _ {h} < 0. 94.
\tag{19.55}
\end{equation*}

这是 $S = \frac{1}{2}$ t--J 模型半经典近似的一次成功。差值 $1 - Z_{h}$ 度量五格点极化子短程关联的量子修正。

\section{$t^{\prime}$--$J$ 模型}

对多于一个空穴，经典基态由填充 $\bar{H}^{t}$ 的 $N^{f}$ 个最低态并极小化 (19.33) 的 $E^{cl}$ 得到。

19.3 节我们发现单极化子不诱导经典 Néel 态的长程畸变；且空穴密度局域于铁磁键附近，极化子之间的经典相互作用是短程的。但必须注意，量子涨落（自旋波）将在 $1/S$ 的更高阶引入更远程的相互作用。

对有限空穴密度，即 $\lim_{\mathcal{N}\to\infty}N^{f}/\mathcal{N}>0$，必须考虑相分离的可能性：体系可以通过把空穴挤进共享 Néel 背景畸变的区域来降低能量\footnote{见 Emery、Kivelson 与 Lin。}。

考虑到原 Hubbard 模型忽略了格点间排斥 Coulomb 相互作用，经典近似得到的短程力与相分离不稳定性在真实铜氧化物层中可能被压制。

然而，凭现有信息可以为极化子准粒子构造有效低能哈密顿量。把 (19.39) 中的 $|\hat{\Omega}^{N}\rangle$ 换成缓变的反铁磁关联组态，计算极化子跳跃率：

\begin{equation*}
\Gamma_{ik} \approx \delta_{s _ {i}, s _ {k}} \left\langle \hat {\Omega} ^ {i} \middle| \bar {\Omega} ^ {i} \right\rangle \bar {\Gamma} _ {ik} \left\langle \bar {\Omega} ^ {k} \middle| \hat {\Omega} ^ {k} \right\rangle ,
\tag{19.56}
\end{equation*}

$\bar{\Gamma}_{ik}$ 是 (19.49) 给出的未微扰隧穿率。$\bar{\Omega}^{i}$ 与 $\bar{\Omega}^{k}$ 表示平移到 $i$、$k$ 为心的五格点极化子，其自旋严格平行于一个共同方向——该方向分别使其与 $\hat{\Omega}^{i}$、$\hat{\Omega}^{k}$ 的交叠最大。(19.56) 中两个交叠因子的相位对每个自旋相消，除非 $\hat{\Omega}^{k}$ 与 $\hat{\Omega}^{i}$ 之间空穴密度不同。用 (19.37)：

\begin{align*}
\langle \hat {\Omega} ^ {i} | \bar {\Omega} ^ {i} \rangle \langle \bar {\Omega} ^ {k} | \hat {\Omega} ^ {k} \rangle & \approx \exp \left[ \frac {i}{2} \sum_ {j} \left( 2S - \rho_{j} ^ {i} \right) \cos \bar {\theta}_{j} ^ {i} \left( \phi_{j} ^ {i} - \bar {\phi}_{j} ^ {i} \right) \right]\\
& \quad \times \exp \left[ - \frac {i}{2} \sum_ {j} \left( 2S - \rho_{j} ^ {k} \right) \cos \bar {\theta}_{j} ^ {k} \left( \phi_{j} ^ {k} - \bar {\phi}_{j} ^ {k} \right) \right]\\
& = \exp \left[ - \frac {i}{2} \cos \bar {\theta} \left( \phi_{i} ^ {i} - \phi_{k} ^ {k} \right) \right] ,
\tag{19.57}
\end{align*}

这里用了一个极化子内总空穴密度为 1 的事实。Néel 规范场 $A^{N}$ 的格子定义已由 (17.11) 给出。这里我们发现极化子跳跃以规范协变方式\footnote{见 Shankar。}耦合到 Néel 规范场：

\begin{equation*}
\Gamma_{ik} = \bar {\Gamma}_{ik} \exp \left( i \eta_{i} \sum_{\mu} A^{N} \cdot \mathbf {x}_{ik} \right) ,
\tag{19.58}
\end{equation*}

$\eta_{i}=\mathrm{sign}(s_{i})$ 是极化子的有效“电荷”，依赖其子晶格指标。

(19.39) 的极化子算符 $p_{i}^{\dagger}$ 在大于两个晶格常量的距离上反对易。于是，长程有序相中稀疏极化子体系的低能性质由

\begin{equation*}
Z \propto \int \mathcal {D} \hat {\Omega} \, \mathcal {D} ^ {2} p \, \exp \left( \int_ {0} ^ {\beta} i \sum_ {i} 2S \mathbf {A} \cdot \dot {\hat {\Omega}} _ {i} - \frac {\bar {J}}{4} \sum_{\langle ij \rangle} \hat {\Omega} _ {i} \cdot \hat {\Omega} _ {j} - L ^ {p} \right)
\tag{19.59}
\end{equation*}

描述，$L^{p}$ 是无自旋 Grassmann 变量 $p_{i}$ 的 Lagrange 量：

\begin{align*}
L ^ {p} & = \sum_ {i} p _ {i} ^ {*} \left( \partial_ {\tau} + i \eta_ {i} A _ {0} ^ {N} + \mu \right) p _ {i} + \sum_{\langle i, j, k \rangle} \Gamma_{ik} e ^ {i \eta_{i} \mathbf {A}^{N} \cdot \mathbf {x}_{ik}} p _ {i} ^ {\dagger} p _ {k}\\
& \quad + \sum_ {ij} U _ {ij} p _ {i} ^ {*} p _ {i} p _ {j} ^ {*} p _ {j} .
\tag{19.60}
\end{align*}

$U_{ij}$ 是双极化子有效相互作用。$L^{p}$ 的规范不变性反映 (19.18) 所示自旋与电荷的局域约束。用满足 (19.10) 的有效从属费米子与 Schwinger 玻色子算符近似极化子与自旋自由度，与 (19.60) 对应的最小哈密顿量是 $t^{\prime}$--$J$ 模型：

\begin{equation*}
\mathcal {H} ^ {t ' - J, s} = P _ {s} \left[ \bar {J} \sum_{\langle ij \rangle} \mathbf {S} _ {i} \cdot \mathbf {S} _ {j} + \sum_{\langle i, k \rangle} \Gamma_{ik} \, p_{i} ^ {\dagger} p_{k} \left( a_{i} ^ {\dagger} a_{k} + b_{i} ^ {\dagger} b_{k} \right) \right] P _ {s} ,
\tag{19.61}
\end{equation*}

自旋 1/2 时：

\begin{equation*}
\mathcal {H} ^ {t ' - J} = P _ {s} \left( \bar {J} \sum_{\langle ij \rangle} \mathbf {S} _ {i} \cdot \mathbf {S} _ {j} - \sum_{\langle i, k \rangle, s} \Gamma_{ik} \, c _ {is} ^ {\dagger} c _ {ks} \right) P _ {s}.
\tag{19.62}
\end{equation*}

$c_{is}$ 是电子算符。$t^{\prime}$--$J$ 模型与初始 t--J 模型的差别在于没有跨子晶格跳跃。$t'$ 跳跃与自旋的局域相互作用较弱，因为空穴在同一子晶格上跳跃时不留下一串反转的自旋。这使得 $t^{\prime}$--$J$ 模型的跳跃项更适于弱耦合处理——至少在反铁磁相中如此。此外，如下文所述，可以用连续近似描述低空穴浓度的低能动力学。

\subsection{超导电性？}

循 Haldane 映射\footnote{见第 12 章。}，自旋相互作用由 $(2+1)$ 维非线性 {$\sigma$} 模型描述，劲度常数与自旋波速度依赖浓度。预期在某个临界密度之上长程自旋序消失，但反铁磁关联在短长度尺度上存留。如 (19.60) 所示，低能自旋涨落通过 Néel 规范场与极化子耦合。Wiegmann、Wen、Shankar 与 Lee（WWSL）为这一耦合体系提出了一个有效场论：

\begin{equation*}
\mathcal {L}^{\mathrm{WWSL}} = \sum_{\alpha = A, B} L^{p, \alpha} [ - i \partial_{\mu} + \eta_{\alpha} A_{\mu}^{N} ] + \frac {1}{m^ {2}} \sum_{\mu \nu} \left( F_{\mu \nu} \right) ^ {2} ,
\tag{19.63}
\end{equation*}

$L^{p,\alpha}$ 是子晶格 $\alpha$ 上极化子的长波 Lagrange 量；手性自旋涨落由 Néel 规范场导出的“电磁场”描述：

\begin{equation*}
F _ {\nu \mu} = \partial_ {\mu} A _ {\nu} ^ {N} - \partial_ {\nu} A _ {\mu} ^ {N}.
\tag{19.64}
\end{equation*}

$m$ 是 Néel 规范场的有效耦合常数，粗略正比于逆自旋关联长度\footnote{见 A. M. Polyakov, \textit{Gauge Fields and Strings} (Harwood, 1987), 8.3 节。}。

WWSL 提出 (19.63) 的基态可能是高温超导体。基本论证是：不同子晶格上的极化子相对 Néel 规范场带相反电荷，在极化子之间产生类电磁吸引。这是一种运动学配对机制，对有效模型的短程相互作用不敏感。而且该方案在磁无序相中依然有效。这尤为令人满意，因为实验上铜氧化物超导体中反铁磁性与超导电性似乎并不共存。

\begin{exercises}

\item 设半径 $R_{c}$ 的大圆形区域内所有自旋铁磁关联。计算大 $\bar{t}/\bar{J}$ 的经典基态能量 $E^{cl}$，对 $R_{c}$ 极小化能量以证明 (19.35)。

\end{exercises}

\begin{chapbib}
\item t--J 模型的 Schwinger 玻色子平均场理论由
  C. Jayaprakash, H. R. Krishnamurthy, and S. Sarker, Phys. Rev. B \textbf{40}, 2610 (1989)；
  C. L. Kane, P. A. Lee, T. K. Ng, B. Chakraborty, and N. Read, Phys. Rev. B \textbf{41}, 2653 (1990)
  完成。
\item 经典极限下极化子大小与相互作用的计算见
  A. Auerbach and B. E. Larson, Phys. Rev. Lett. \textbf{66}, 2262 (1991)。
\item 空穴在涨落反铁磁背景中的规范耦合，最早在连续表述中导出于
  R. Shankar, Nucl. Phys. B \textbf{330}, 433 (1990)。
\item 用受限 Green 函数（“路径分解展开”）计算多维隧穿率的综述：
  A. Auerbach and S. Kivelson, Nucl. Phys. Rev. B \textbf{257}, 799 (1985)。
\item 极化子算符的定义及其权重在小格子上对 t--J 模型的数值计算：
  E. Dagotto and J. R. Schrieffer, Phys. Rev. B \textbf{43}, 8705 (1991)。
\item t--J 模型相分离的变分论证：
  V. J. Emery, S. A. Kivelson, and H. Q. Lin, Phys. Rev. Lett. \textbf{64}, 475 (1990)。
\item $t^{\prime}$--$J$ 模型高温超导理论由
  P. B. Wiegmann, Phys. Rev. Lett. \textbf{60}, 821 (1988)；
  P. A. Lee, Phys. Rev. Lett. \textbf{63}, 690 (1989)；
  X.-G. Wen, Phys. Rev. B \textbf{39}, 7223 (1989)
  提出。
\item 二维反铁磁体中空穴问题的其他途径：
  S. Trugman, Phys. Rev. B \textbf{41}, 892 (1990)；
  B. Schraiman and E. Siggia, Phys. Rev. Lett. \textbf{62}, 156 (1989)。
\end{chapbib}
